We introduced a prefix-level oracle benchmark for verifiable reasoning and used it to test a simple but consequential hypothesis:
is prefix compute allocation well modeled as ordered scalar thresholding?
Our answer is no at the benchmark level and conditional at the controller-ranking level.

The benchmark shows substantial high-determinacy action crossing in two exact-checker domains.
That is enough to reject simple continuation-score thresholding, and more generally a single ordered confidence score, as a sufficient description of low-regret prefix control.
At the same time, the strongest structured-controller result is a deployment-gap diagnosis rather than a universal win:
exact-state structure is clearly valuable, but predicted-state performance degrades sharply under noisy state estimation.

This leads to the paper's main message.
Prefix compute allocation in verifiable reasoning should be studied as \processstate control.
The right benchmark is prefix-level and oracle-based.
The right evaluation target is budgeted decision quality, not only final-answer accuracy.
And the main open systems problem is reliable state identification under deployment noise.

The broader takeaway is methodological.
If we evaluate only final-answer accuracy, many of these distinctions disappear.
If we evaluate prefix actions against oracle regret, revision harm, and compute-value calibration, they become explicit.
That is the role of the benchmark in this paper.
It does not merely rank controllers; it separates mechanism from deployment, and action quality from raw confidence, which is exactly what scalar-threshold narratives tend to collapse.
